\documentclass[11pt,a4paper,sans]{moderncv}
%\pdfminorversion=5
%\pdfobjcompresslevel=3 
%\pdfcompresslevel=9

\usepackage{fontspec}
\setmainfont{Raleway}

\usepackage{dtk-logos}
% moderncv themes
\moderncvstyle{casual}   
\moderncvcolor{black}                
% adjust the page margins
\usepackage[scale=0.75]{geometry}
\setlength{\hintscolumnwidth}{3.6cm}
\recomputelengths
%\setlength{\hintscolumnwidth}{3cm}                % if you want to change the width of the column with the dates
%\setlength{\makecvtitlenamewidth}{10cm}           % for the 'classic' style, if you want to force the width allocated to your name and avoid line breaks. be careful though, the length is normally calculated to avoid any overlap with your personal info; use this at your own typographical risks...

% personal data
\name{Dr. Patrick}{Diehl}
\title{Curriculum Vit\ae{}}                               % optional, remove / comment the line if not wanted
\address{Digital Media Center, 340 E Parker Blvd}{Baton Rouge, LA 70803, USA}{}% optional, remove / comment the line if not wanted; the "postcode city" and "country" arguments can be omitted or provided empty
%\phone[mobile]{+49~(176)~32311017}                   % optional, remove / comment the line if not wanted; the optional "type" of the phone can be "mobile" (default), "fixed" or "fax"
\phone[fixed]{+1~(225)~578~5061}
%\phone[fax]{+49~(228)~73~3979}
\email{patrickdiehl@lsu.edu}    
 % optional, remove / comment the line if not wanted
\homepage{www.diehlpk.de}                         % optional, remove / comment the line if not wanted
\social[linkedin]{patrickdiehl}                        % optional, remove / comment the line if not wanted
\social[twitter]{diehlpk}                             % optional, remove / comment the line if not wanted
\social[github]{diehlpk}                              % optional, remove / comment the line if not wanted
%\social[skype]{diehlpk}
\social[orcid]{0000-0003-3922-8419}
%\extrainfo{additional information}                 % optional, remove / comment the line if not wanted
\photo[90pt][0.4pt]{picture}                       % optional, remove / comment the line if not wanted; '64pt' is the height the picture must be resized to, 0.4pt is the thickness of the frame around it (put it to 0pt for no frame) and 'picture' is the name of the picture file
%\quote{Some quote}                                 % optional, remove / comment the line if not wanted


% bibliography with mutiple entries
\usepackage{multibib}
\newcites{proceedings,journal,talk,thesis,reports,exp,pre,spapers,booksedited}{{\underline{Series- and conference contributions}},{\underline{Journal articles}},{},{},{\underline{Technical reports}},{},{\underline{Preprints}},{\underline{Short papers}},{\underline{Edited books}}}

\usepackage{enumitem}

%----------------------------------------------------------------------------------
%            content
%----------------------------------------------------------------------------------
\begin{document}

\makecvtitle

%%%%%%%%%%%%%%%%%%%%%%%%%%%%%%%%%%%%%%%%%%%%%%%%%%%%%%%%%%%%%%%%%%%%%%%%%%%%%%
\section{Education}
%%%%%%%%%%%%%%%%%%%%%%%%%%%%%%%%%%%%%%%%%%%%%%%%%%%%%%%%%%%%%%%%%%%%%%%%%%%%%%
\cventry{2017}{PhD}{Applied mathematics}{University of Bonn}{Germany}{}
\cventry{2012}{Diploma}{Computer Science}{University of Stuttgart}{Germany}{}
%%%%%%%%%%%%%%%%%%%%%%%%%%%%%%%%%%%%%%%%%%%%%%%%%%%%%%%%%%%%%%%%%%%%%%%%%%%%%%
\section{Awards and Honors}
%%%%%%%%%%%%%%%%%%%%%%%%%%%%%%%%%%%%%%%%%%%%%%%%%%%%%%%%%%%%%%%%%%%%%%%%%%%%%%
\cvitem{2019}{IEEE SCIVIS Contest 2019, First Prize, Visual Analysis of Structure Formation in Cosmic Evolution, \href{https://www.youtube.com/watch?v=ykn3ewqWUcw}{Video}, \href{https://megamol.org/wp-content/uploads/2019/10/poster2-reduced.jpg}{Poster}, and \href{https://ieeexplore.ieee.org/abstract/document/8968855}{Short paper}}
%%%%%%%%%%%%%%%%%%%%%%%%%%%%%%%%%%%%%%%%%%%%%%%%%%%%%%%%%%%%%%%%%%%%%%%%%%%%%%
\section{Grant history}
%%%%%%%%%%%%%%%%%%%%%%%%%%%%%%%%%%%%%%%%%%%%%%%%%%%%%%%%%%%%%%%%%%%%%%%%%%%%%%

%%%%%%%%%%%%%%%%%%%%%%%%%%%%%%%%%%%%%%%%%%%%%%%%%%%%%%%%%%%%%%%%%%%%%%%%%%%%%%
\subsection{Current Research (chronological order; most recent one first)}
%%%%%%%%%%%%%%%%%%%%%%%%%%%%%%%%%%%%%%%%%%%%%%%%%%%%%%%%%%%%%%%%%%%%%%%%%%%%%%
\cvitem{}{
\begin{enumerate}
\item
\begin{description}
\item Grant \#2229751 (Rod Tohid)
\item Name of Funding Organization: National Science Foundation
\item Amount Awarded: \$300,000
\item Period of Grant Award: Sept 15 2022 - Oct 31 2023
\item Title of Project: POSE: Phase I: Constellation: A Pathway to Establish the STE$\vert\vert$AR Open-Source Organization
\item Role on Project: Co-PI
\end{description}
\end{enumerate}
}

%%%%%%%%%%%%%%%%%%%%%%%%%%%%%%%%%%%%%%%%%%%%%%%%%%%%%%%%%%%%%%%%%%%%%%%%%%%%%%
\cvitem{}{Completed Research (chronological order; most recent one first)}
%%%%%%%%%%%%%%%%%%%%%%%%%%%%%%%%%%%%%%%%%%%%%%%%%%%%%%%%%%%%%%%%%%%%%%%%%%%%%%
\begin{enumerate}[leftmargin=4.5cm]
\item Grant \#524125 (Hartmut Kaiser)
\begin{description}
\item Name of Funding Organization: Pacific Northwest National Laboratory
\item Amount Awarded: \$50,000
\item Period of Grant Award: June 25 - Oct 31 2020
\item Title of Project: High Performance Data Analytics (HPDA) Scalable Second-Order Optimization (SSO)
\item Role on Project: Co-PI
\end{description}
\end{enumerate}
%%%%%%%%%%%%%%%%%%%%%%%%%%%%%%%%%%%%%%%%%%%%%%%%%%%%%%%%%%%%%%%%%%%%%%%%%%%%%%
\section{Allocations}
%%%%%%%%%%%%%%%%%%%%%%%%%%%%%%%%%%%%%%%%%%%%%%%%%%%%%%%%%%%%%%%%%%%%%%%%%%%%%%

%%%%%%%%%%%%%%%%%%%%%%%%%%%%%%%%%%%%%%%%%%%%%%%%%%%%%%%%%%%%%%%%%%%%%%%%%%%%%%
\cvitem{}{Current Allocations (chronological order; most recent one first)}
%%%%%%%%%%%%%%%%%%%%%%%%%%%%%%%%%%%%%%%%%%%%%%%%%%%%%%%%%%%%%%%%%%%%%%%%%%%%%%
\begin{enumerate}[leftmargin=4.5cm]
\item  Project xpress (Alice Koniges)
\begin{description}
\item Project xpress (Alice Koniges)
\item Title of Project: HPX and OpenMP
\item Amount awarded:  CPU node hours 4300 and GPU node hours 1700
\item Cluster: Perlmutter @ NERSC
\item Role on Project: Senior personal
\end{description}
\end{enumerate}


%%%%%%%%%%%%%%%%%%%%%%%%%%%%%%%%%%%%%%%%%%%%%%%%%%%%%%%%%%%%%%%%%%%%%%%%%%%%%%
\cvitem{}{Completed Allocations (chronological order; most recent one first)}
%%%%%%%%%%%%%%%%%%%%%%%%%%%%%%%%%%%%%%%%%%%%%%%%%%%%%%%%%%%%%%%%%%%%%%%%%%%%%%

\begin{enumerate}[leftmargin=4.5cm]
\item  Project hp210311 (Patrick Diehl)
\begin{description}
\item Title of Project: Porting Octo-Tiger, an astrophysics program simulating the evolution of star systems based on the fast multipole method on adaptive Octrees
\item Type: Test-bed
\item Amount awarded: 21k node hours
\item Cluster: Fugaku @ RIKEN Center for Computational Science, Japan
\item Role on Project: PI
\end{description}
\item  Project PaDi032321F (Patrick Diehl)
\begin{description}
\item Title of Project: Porting Octo-Tiger, an astrophysics program simulating the evolution of star systems based on the fast multipole method on adaptive Octrees
\item Type: Test-bed
\item Amount awarded: 10k node hours
\item Cluster: Ookami @ Stonybrook  University,  USA
\item Role on Project: PI
\end{description}
\end{enumerate}

%%%%%%%%%%%%%%%%%%%%%%%%%%%%%%%%%%%%%%%%%%%%%%%%%%%%%%%%%%%%%%%%%%%%%%%%%%%%%%
\section{Editorial duties}
%%%%%%%%%%%%%%%%%%%%%%%%%%%%%%%%%%%%%%%%%%%%%%%%%%%%%%%%%%%%%%%%%%%%%%%%%%%%%%
\cventry{06/20--current}{Topic editor}{Computational fracture mechanics, Applied mathematics, C\texttt{++}, asynchronous and task-based programming}{The Journal of Open Source Software}{}{}{}
\cventry{2022}{Guest editor}{Special issue: Science Gateways: Accelerating Research and Education, Computing in Science \& Engineering}{}{}{}
\cventry{2021}{Guest editor}{Special issue: Peridynamics and its Current Progress,  Computer Modeling in Engineering \& Sciences (CMES)}{}{}{}
%%%%%%%%%%%%%%%%%%%%%%%%%%%%%%%%%%%%%%%%%%%%%%%%%%%%%%%%%%%%%%%%%%%%%%%%%%%%%%
\section{Research experience}
%%%%%%%%%%%%%%%%%%%%%%%%%%%%%%%%%%%%%%%%%%%%%%%%%%%%%%%%%%%%%%%%%%%%%%%%%%%%%%
\cventry{05/22--current}{Adjunct faculty}{Department of Physics \& Astronomy}{Louisiana State University}{Baton Rouge, LA, USA}{}
\cventry{10/18--current}{Research scientist}{Center for Computation \& Technology}{Louisiana State University}{Baton Rouge, LA, USA}{}
\cventry{02/17--09/18}{Postdoctoral fellow}{Laboratory of Multiscale Mechanics}{Polytechnique Montr\'eal}{QC, Canada}{}
\cvitem{}{\begin{itemize}
\item Benchmark peridynamic simulations against experimental data for composite materials
\item Extracting constitutive mechanical parameters in linear elasticity using the virtual fields method within the ordinary state-based peridynamic framework (with Rolland Delorme)
\item  Hybrid image processing approach for crack area detection and tracking using local Digital Image Correlation results (with Ilyass Tabiai)
\end{itemize}}
\cventry{04/13--01/17}{Research Assistant}{Institute for Numerical Simulation}{University Bonn}{Bonn, Germany}{}
\cvitem{}{\begin{itemize}
\item Modeling and simulation of crack and fractures in solids using peridynamic
\end{itemize}}
\cventry{07/12--03/13}{Research Assistant}{Institute for Simulation of large Systems}{University Stuttgart}{Stuttgart, Germany}{}
\cventry{10/07--06/12}{Student}{Computer Science (Major software engineering)}{University Stuttgart}{Stuttgart, Germany}{}
%%%%%%%%%%%%%%%%%%%%%%%%%%%%%%%%%%%%%%%%%%%%%%%%%%%%%%%%%%%%%%%%%%%%%%%%%%%%%%
\section{Visiting positions}
%%%%%%%%%%%%%%%%%%%%%%%%%%%%%%%%%%%%%%%%%%%%%%%%%%%%%%%%%%%%%%%%%%%%%%%%%%%%%%
\cventry{2015}{Guest Research Assistant}{Center for Computation and Technology}{Louisiana State University}{}{}
%%%%%%%%%%%%%%%%%%%%%%%%%%%%%%%%%%%%%%%%%%%%%%%%%%%%%%%%%%%%%%%%%%%%%%%%%%%%%%
\section{Research Interests}
%%%%%%%%%%%%%%%%%%%%%%%%%%%%%%%%%%%%%%%%%%%%%%%%%%%%%%%%%%%%%%%%%%%%%%%%%%%%%%
\cvitem{Computational engineering}{}

\begin{itemize}[leftmargin=4cm]
\item Peridynamics theory for the application in solids, like glassy or composite materials,
\item Validation and verification of simulations against experimental data,
\item Assembly of experimental data for comparison with simulations,
\item Application of machine learning to experiments and simulations.
\end{itemize}
\cvitem{High Performance Computing}{}
\begin{itemize}[leftmargin=4cm]
\item The C\texttt{++} Standard Library for Parallelism and Concurrency (HPX),
\item Asynchronous many task systems and there application in computational engineering.
\item Efficient, performance portable, and scalable High-Performance Parallel Programming using Modern C++.
\end{itemize}
\cvitem{Open science}{}
\begin{itemize}[leftmargin=4cm]
\item Open Source Software for scientific applications,
\item Open data for sharing raw experimental results.
\end{itemize}
%%%%%%%%%%%%%%%%%%%%%%%%%%%%%%%%%%%%%%%%%%%%%%%%%%%%%%%%%%%%%%%%%%%%%%%%%%%%%%
\section{Teaching experience}
%%%%%%%%%%%%%%%%%%%%%%%%%%%%%%%%%%%%%%%%%%%%%%%%%%%%%%%%%%%%%%%%%%%%%%%%%%%%%%
\cvitem{Instructor}{}
\begin{itemize}[leftmargin=4cm]
\item Research Technologies and Methods (MEDP 7098), Department of Physics \& Astronomy, Louisiana State University, Taught: 2022
\item Parallel computational mathematics (M 4997), Department of Mathematics,  Louisiana State University, Taught: 2019, 2020, and 2021
\end{itemize}

\cvitem{Teaching assistant}{}
\begin{itemize}[leftmargin=4cm]
\item Einf\"uhrung in die Numerische Mathematik (Introduction to numerical mathematics), University of Bonn, 2015
\item Algorithmische Mathematik (Mathematical algorithms), University of Bonn, 2013/2014
\item Wissenschaftliches Rechnen 2 (Scientific Computing 2),  University of Bonn, 2013 
\end{itemize}


%%%%%%%%%%%%%%%%%%%%%%%%%%%%%%%%%%%%%%%%%%%%%%%%%%%%%%%%%%%%%%%%%%%%%%%%%%%%%%
\section{Certificates}
%%%%%%%%%%%%%%%%%%%%%%%%%%%%%%%%%%%%%%%%%%%%%%%%%%%%%%%%%%%%%%%%%%%%%%%%%%%%%%
\cvitem{Baden--W\"urttemberg Certificate}{Certificate for successful completion of the program in higher education pedagogy by the center for educational development of the state of Baden--W\"urttemberg.}

%%%%%%%%%%%%%%%%%%%%%%%%%%%%%%%%%%%%%%%%%%%%%%%%%%%%%%%%%%%%%%%%%%%%%%%%%%%%%%
\section{Advising and related student services}
%%%%%%%%%%%%%%%%%%%%%%%%%%%%%%%%%%%%%%%%%%%%%%%%%%%%%%%%%%%%%%%%%%%%%%%%%%%%%%
\cvitem{Co-supervised theses}{}
\begin{itemize}[leftmargin=4cm]
\item University of Stuttgart:
\begin{itemize}
\item Pfander, David: Eine künstliche Intelligenz für das Kartenspiel Tichu, Studienarbeit Nr. 2398, 2013.
\item Kanis, Sebastian: GPU-based Numerical Integration in the Partition of Unity Method, Diplomarbeit Nr. 3405, 2013.
\end{itemize}
\end{itemize}

\cvitem{Graduate Committee Member}{}
\begin{itemize}[leftmargin=4cm]
\item Master thesis: M. Reeser (LSU, 2020)
\item Honors project: J. Trepper (LSU, 2020)
\end{itemize}

\cvitem{Google Summer of Code}{}
\begin{itemize}[leftmargin=4cm]
\item University of Bonn: A. Nigam (2016)
\item Polytechnique Montreal:  M. Seshadri (2017), G. Laberge and J. Golinowski (2018)
\item Louisiana State University: P. Gadika (2020)
\end{itemize}

\cvitem{Google Season of Docs}{}
\begin{itemize}[leftmargin=4cm]
\item Louisiana State University: R. Stobaugh (2019)
\end{itemize}

\cvitem{NSF REU}{}
\begin{itemize}[leftmargin=4cm]
\item Louisiana State University: A. Edwards (2021) and E. Downing (2022)
\end{itemize}


%%%%%%%%%%%%%%%%%%%%%%%%%%%%%%%%%%%%%%%%%%%%%%%%%%%%%%%%%%%%%%%%%%%%%%%%%%%%%%
\section{Academic-related Professional and Public Service}
%%%%%%%%%%%%%%%%%%%%%%%%%%%%%%%%%%%%%%%%%%%%%%%%%%%%%%%%%%%%%%%%%%%%%%%%%%%%%%
\cvitem{10/21--current}{Reviewer for Mathematical Reviews\textsuperscript{\textregistered} (MathSciNet\textsuperscript{\textregistered})}
\cvitem{07/21--current}{Member-at-Large of the USACM Technical Thrust Area (TTA) on Large Scale Structural Systems and Optimal Design (Large-Scale)}

\cvitem{03/20--current}{Liaison for the Louisiana district of the SIAM Texas-Louisiana Section\linebreak
Duties:}
\begin{itemize}[leftmargin=4cm]
\item Making sure that people at universities, research institutions and industry in your district know
  about our activities and getting their suggestions on what we can do better
\item Serving on the organizing committee for the annual meeting
\end{itemize}
\cvitem{10/17--09/18}{ASSEP Labor relations officers for postdoctoral fellows}

%%%%%%%%%%%%%%%%%%%%%%%%%%%%%%%%%%%%%%%%%%%%%%%%%%%%%%%%%%%%%%%%%%%%%%%%%%%%%%
\section{Organization of Conferences, Workshops, and Symposia}
%%%%%%%%%%%%%%%%%%%%%%%%%%%%%%%%%%%%%%%%%%%%%%%%%%%%%%%%%%%%%%%%%%%%%%%%%%%%%%
\cvitem{Committee}{}
\begin{itemize}[leftmargin=4cm]
\item Science Gateways Conference 2021, Publication chair
\item The International Conference for High Performance Computing, Networking, Storage, and Analysis (SC) 21, AD/AE Appendices
\item Science Gateways Conference 2021, Program committee member
\item The International Conference for High Performance Computing, Networking, Storage, and Analysis (SC) 21, Virtual logistics
\end{itemize}

\cvitem{Symposia}{}
\begin{itemize}[leftmargin=4cm]
\item Modeling and Simulation for Complex Material Behavior, 14th U.S. National Congress on Computational Mechanics, \href{http://14.usnccm.org/MS402}{Link}.
\item Peridynamic Theory and Multiscale Methods for Complex Material Behavior, 14th World Congress on Computational Mechanics (WCCM XIV).
\item Peridynamic Theory and Multiscale Methods for Complex Material Behavior, 16th National Congress on Computational Mechanics. 
\item Nonlocal Models in Mathematics and Computation, 3rd Annual Meeting of the SIAM Texas-Louisiana Section, 2020 
\item Peridynamic Theory and Multiscale Methods for Complex Material Behavior, 15th World Congress on Computational Mechanics (WCCM XV).
\item Recent Developments in Peridynamics Modeling, 19th U.S. National Congress on Theoretical and Applied Mechanics.
\item  Peridynamic Theory and Multiscale Methods for Complex Material Behavior,  9th GACM Colloquium on Computational Mechanics
\item Theoretical and Computational Aspects of Nonlocal Operator, 7th Annual Meeting of SIAM Central States Section
\item Nonlocal Models in Mathematics and Computation, 5th Annual Meeting of the SIAM Texas-Louisiana Section, 2022
\item  Recent Developments in Peridynamics Modeling, 17th U.S. National Congress on Computational Mechanics, 2023
\item Nonlocal Modeling, Analysis, and Computation, 10th International Congress on Industrial and Applied Mathematics, 2023
\end{itemize}

\cvitem{Workshop}{}
\begin{itemize}[leftmargin=4cm]
\item Workshop on Experimental and Computational Fracture Mechanics:  Validating peridynamics and phase field models for fracture
prediction and experimental design, \href{http://wfm2020.usacm.org/}{Link}. Sponsored by
\begin{itemize}
\item US Association for Computational Mechanics,
\item Center for Computation \& Technology at Louisiana State University, 
\item Oak Ridge National Laboratory,
\item Society for Experimental Mechanics,
\item U.S. National Committee on Theoretical and Applied Mechanics (AmeriMech)
\end{itemize}
\item Asynchronous Many-Task systems for Exascale 2021 held in conjunction with \href{https://2021.euro-par.org/}{Euro-Par 2021},  \href{https://amte2021.stellar-group.org}{Link}
\item Asynchronous Many-Task systems for Exascale 2022 held in conjunction with \href{https://2022.euro-par.org/}{Euro-Par 2022},  \href{https://amte2022.stellar-group.org}{Link}
\item Workshop on Workshop on Asynchronous Many-Task Systems and Applications 2023,  \href{https://wamta23.stellar-group.org/}{Link}
\item Asynchronous Many-Task systems for Exascale 2023 held in conjunction with \href{https://2023.euro-par.org/}{Euro-Par 2023},  \href{https://amte2023.stellar-group.org}{Link}
\end{itemize}

\cvitem{Panel}{}
\begin{itemize}[leftmargin=4cm]
\item Joint USACM Large-Scale TTA – EMI Computational Mechanics Committee Career Path Panel, Speaker, Engineering Mechanics Institute Conference 2022
\item TBAA: Task-Based Algorithms and Applications, Moderator, "International Conference for High Performance Computing, Networking, Storage and Analysis (SC)" 2020. \href{https://stellar-group.org/research/tbaa2020-scpanel/}{Link}
\item AI Ethics/Algorithmic Justice, Organizer, Colloquium on Artificial Intelligence Research and Optimization,  Louisiana State University. 
\end{itemize}


\cvitem{Meeting}{}
\begin{itemize}[leftmargin=4cm]
\item 3rd Annual Meeting of the SIAM Texas-Louisiana Section, October 16 - 18,  2020. \href{https://www.math.tamu.edu/conferences/SIAMTXLA/}{Link}.
\item 4th Annual Meeting of the SIAM Texas-Louisiana Section,  November,  5 - 7,  2021 \href{https://faculty.utrgv.edu/eleftherios.gkioulekas/2021-siam-txla-meeting-spi/index.html}{Link}.
\item 5th Annual Meeting of the SIAM Texas-Louisiana Section,  November,  4 - 6,  2022 \href{https://www.math.uh.edu/siamtxla22/index.shtml}{Link}.
\end{itemize}

\cvitem{Colloquium}{}
\begin{itemize}[leftmargin=4cm]
\item Colloquium on Artificial Intelligence Research and Optimization, Louisiana State University. \href{https://stellar-group.org/research/distributed-machine-learning/}{Link}.
\item Large-Scale TTA Early-Career Colloquium, USACM. \href{https://unsacm.memberclicks.net/tta--large-scale-structural-systems }{Link}
\end{itemize}

\cvitem{Short course}{}
\begin{itemize}[leftmargin=4cm]
\item SC16-001 Advanced Parallel Programming in C\texttt{++}, 16th U.S. National Congress on Computational Mechanics
\item Advanced Parallel Programming in C\texttt{++}, 15th World Congress on Computational Mechanics
\end{itemize}

\cvitem{Mentoring events}{}
\begin{itemize}[leftmargin=4cm]
\item 15th World Congress on Computational Mechanics, 2022
\item 5th Annual Meeting of the SIAM Texas-Louisiana Section, 2022
\end{itemize}

%%%%%%%%%%%%%%%%%%%%%%%%%%%%%%%%%%%%%%%%%%%%%%%%%%%%%%%%%%%%%%%%%%%%%%%%%%%%%%
\section{Conference and Workshop Grants}
%%%%%%%%%%%%%%%%%%%%%%%%%%%%%%%%%%%%%%%%%%%%%%%%%%%%%%%%%%%%%%%%%%%%%%%%%%%%%%
\cvitem{2020}{\textbf{AmeriMech symposium}: Experimental and Computational Fracture Mechanics: Validating peridynamics and phase field models for fracture prediction and experimental design (\$4000)}

%%%%%%%%%%%%%%%%%%%%%%%%%%%%%%%%%%%%%%%%%%%%%%%%%%%%%%%%%%%%%%%%%%%%%%%%%%%%%%
\section{Publications}
%%%%%%%%%%%%%%%%%%%%%%%%%%%%%%%%%%%%%%%%%%%%%%%%%%%%%%%%%%%%%%%%%%%%%%%%%%%%%%
\nocitepre{marcello2021octotiger,diehl2021performance}

\bibliographystylepre{../list/habbrvyr}
\bibliographypre{../list/publications}     
%%%%%%%%%%%%%%%%%%%%%%%%%%%%%%%%%%%%%%%%%%%
\nocitebooksedited{europar2021,europar2022,europar2023a,europar2023b,wamta23,wamta24,10197475,10231033,usrse2025,wamta25,europar25a,europar25b}

\bibliographystylebooksedited{../list/habbrvyr}
\bibliographybooksedited{../list/publications}     
%%%%%%%%%%%%%%%%%%%%%%%%%%%%%%%%%%%%%%%%%%%
\nocitejournal{diehl2021octo,10.1093/mnras/stab937,Diehl2016FRAC,busler2017,diehl2018long,Diehl2019survey,TABIAI2019106485,DaiB:2019:PDS:3295500.3356221,Delorme2020,Kaiser2020,PRUDHOMME2020113391,diehl2018implementation,9405442,Jha2021,https://doi.org/10.1002/cpe.6893} 

\bibliographystylejournal{../list/habbrvyr}
\bibliographyjournal{../list/publications}  
\nocitebook{Diehl.Schweitzer:2014,Franzelin.Diehl.Pfluger:2014,Diehl.Schweitzer:2015,heller2016,Diehl2017a,Heller2017,diehl2018integration,tohid2018asynchronous,Zhang:2019:IHM:3318170.3318191}
\bibliographystyleproceedings{../list/habbrvyr}
\bibliographyproceedings{../list/publications} 
\nocitespapers{diehl2020gateways,schatz2019visual,10.1115/IMECE2021-70793} 
\bibliographystylespapers{../list/habbrvyr}
\bibliographyspapers{../list/publications} 
\nocitereports{Diehl.Lipton.Schweitzer:2016,zhang2020supporting,wfm2020-report,tbaa20-report}
    

\bibliographystylereports{../list/habbrvyr}
\bibliographyreports{../list/publications}          
%%%%%%%%%%%%%%%%%%%%%%%%%%%%%%%%%%%
\section{Invited talks and Presentations}
\nocitetalk{Diehl:2014,Diehl:2014*3,Diehl:2013,Diehl:2014*1,Diehl:2014*2,Diehl:2013*2,Diehl:2013*1,Diehl:2015*2,Diehl:2016SIAM,Diehl:2016-mafleap,Diehl:2016-wccm,Diehl:2015,Diehl:2016-sandiego,Diehl:2017-him,Diehl2017MontrealCanada,Diehl2017TexasAustin,Diehl:2018*2,Diehl:2018*3,Diehl:2018-wccm,Diehl:2018-sc,Diehl:2019-emi}
\bibliographystyletalk{../list/habbrvyr}
\bibliographytalk{../list/publications}
%%%%%%%%%%%%%%%%%%%%%%%%%%%%%%%%%%%%%%%%%%%%%%%%%%%%%%%%%%%%%%%%%%%%%%%%%%%%%%
\section{Thesis}
\nocitethesis{Diehl2017,Diehl:2012}
\bibliographystylethesis{../list/habbrvyr}
\bibliographythesis{../list/publications} 
%%%%%%%%%%%%%%%%%%%%%%%%%%%%%%%%%%%%%%%%%%%%%%%%%%%%%%%%%%%%%%%%%%%%%%%%%%%%%%
\section{Open source software}
\cvitem{}{}
\begin{itemize}[leftmargin=4cm]
\item NLMech - Implementation of finite element and finite difference approximation of Nonlocal models \href{https://doi.org/10.21105/joss.03020}{\includegraphics[scale=0.5]{status-nlmech-small.pdf}}
\item HPX - C\texttt{++} standard library for parallelism and concurrency \href{https://doi.org/10.21105/joss.02352}{\includegraphics[scale=0.5]{status-hpx.pdf}}
\item Octo-Tiger - GPU-accelerated astrophysics code to simulate the merging of stars
\end{itemize}


\section{Raw experimental data}
\nociteexp{ilyass_tabiai_2018_1172068}
\bibliographystyleexp{../list/habbrvyr}
\bibliographyexp{../list/publications} 
%%%%%%%%%%%%%%%%%%%%%%%%%%%%%%%%%%%%%%%%%%%%%%%%%%%%%%%%%%%%%%%%%%%%%%%%%%%%%%
\section{Professional Organizations}
%%%%%%%%%%%%%%%%%%%%%%%%%%%%%%%%%%%%%%%%%%%%%%%%%%%%%%%%%%%%%%%%%%%%%%%%%%%%%%
\cvitem{}{}
\begin{itemize}[leftmargin=4cm]
\item Society for Industrial and Applied Mathematics (SIAM)
\item U.S. Association for Computational Mechanics (USACM)
\item Informatik-Forum Stuttgart e. V.
\end{itemize}

%%%%%%%%%%%%%%%%%%%%%%%%%%%%%%%%%%%%%%%%%%%%%%%%%%%%%%%%%%%%%%%%%%%%%%%%%%%%%%
\section{Reviewer}
%%%%%%%%%%%%%%%%%%%%%%%%%%%%%%%%%%%%%%%%%%%%%%%%%%%%%%%%%%%%%%%%%%%%%%%%%%%%%%
\cvitem{}{International Journal of Mechanical Sciences, Fatigue \& Fracture of Engineering Materials \& Structures, Computer Physics Communications, International Journal of Fracture, Parallel Computing, International Journal of High Performance Computing Applications, Computer Methods in Applied Mechanics and Engineering, Theoretical and Applied Fracture Mechanics,   Mathematics and Mechanics of Solids,  Journal of Engineering Mechanics, Computational Mechanics,  Mathematical Reviews,  Engineering with Computers,  Computer Modeling in Engineering \& Sciences (CMES),  Journal of Peridynamics and Nonlocal Modeling,  Engineering Analysis with Boundary Elements,  and Concurrency and Computation: Practice and Experience.}
%%%%%%%%%%%%%%%%%%%%%%%%%%%%%%%%%%%%%%%%%%%%%%%%%%%%%%%%%%%%%%%%%%%%%%%%%%%%%%
\section{References}
%%%%%%%%%%%%%%%%%%%%%%%%%%%%%%%%%%%%%%%%%%%%%%%%%%%%%%%%%%%%%%%%%%%%%%%%%%%%%%
\cvitem{Phd advisers}{}
\begin{itemize}[leftmargin=4cm]
\item Dr.~Marc~Alexander~Schweitzer, Institute for Numerical Simulation, University of Bonn, Germany, E-Mail: \url{schweitzer@ins.uni-bonn.de}
\item Dr.~Daniel~Peterseim, Numerische Mathematik, University of Augsburg, Germany, E-Mail: \url{daniel.peterseim@math.uni-augsburg.de}
\end{itemize}
\cvitem{Postdoctoral fellow adviser}{}
\begin{itemize}[leftmargin=4cm]
\item Dr.~Serge~Prudhomme, Department of Mathematical and Industrial Engineering, Polytechnique Montr\'eal, Canada, E-Mail: \url{serge.prudhomme@polymtl.ca}
\end{itemize}
\cvitem{Collaborators}{}
\begin{itemize}[leftmargin=4cm]
\item Dr.~Robert~Lipton, Mathematics of Materials Science, Louisiana State University, USA, E-Mail: \url{lipton@math.lsu.edu}
\item Dr.~Thomas~Ertl, Visualization and Interactive Systems Institute, University of Stuttgart, Germany, E-Mail: \url{thomas.ertl@vis.uni-stuttgart.de}
\item Dr.~Hartmut~Kaiser, Department of Computer Science, Louisiana State University, USA, E-Mail: \url{hkaiser@cct.lsu.edu}
\item Dr.~Juhan~Frank, Department of Physics \& Astronomy, Louisiana State University, USA, E-Mail: \url{frank@rouge.phys.lsu.edu}
\item Dr.~Steven~R.~Brandt, Center of Computation \& Technology, Louisiana State University, USA, E-Mail: \url{sbrandt@cct.lsu.edu}
\item Dr.~Ilyass~Tabiai, École de Technologie Supérieur, Canada, \\E-Mail: \url{ilyass.tabiai@etsmtl.net}
\item Dr.~Steffen~Frey, University of Groningen, Netherlands, E-Mail: \url{s.d.frey@rug.nl}
\item Dr.~Dirk~Pfl\"uger, Institute for Parallel and Distributed Systems, University of Stuttgart, Germany, Netherlands, E-Mail: \url{Dirk.Pflueger@ipvs.uni-stuttgart.de}
\item Dr.~Mayank Tagi, Craft \& Hawkins Department of Petroleum Engineering, Louisiana State University, USA, E-Mail:  \url{mtyagi@lsu.edu}
\item Dr.~Geoffrey C. Clayton, Department of Physics and Astronomy, Louisiana State University, USA, E-Mail: \url{gclayton@fenway.phys.lsu.edu}
\item Dr.~Orsola De Marco, Department of Physics and Astronomy, Macquarie University, Australia, E-Mail: \url{orsola.demarco@mq.edu.au}
\item Dr.~Thomas Wick, Leibniz Universit{\"a}t Hannover, Institut f{\"u}r Angewandte Mathematik, Germany, E-Mail: \url{thomas.wick@ifam.uni-hannover.de}
\end{itemize}
\begin{center}
Last update: \today
\end{center}
\end{document}


%% end of file `template.tex'.
